\chapter{Introdução}

A Inteligência Artificial (IA) alcançou avanços expressivos nos últimos anos, impulsionada principalmente pelo aprendizado profundo (\textit{deep learning}) e técnicas baseadas em redes neurais artificiais \cite{young_recent_2018, otter_survey_2021}. Esses métodos têm superado abordagens tradicionais em diversas aplicações, especialmente em tarefas complexas relacionadas à decisão, classificação e extração de informações.

Particularmente no Processamento de Linguagem Natural (PLN), esses avanços transformaram profundamente o cenário tecnológico \cite{torfi_natural_2021}. Modelos neurais profundos, como as redes neurais recorrentes (RNNs), redes LSTM (\textit{Long Short-Term Memory}) e mecanismos de atenção, mostraram-se altamente eficazes em tarefas como tradução automática, análise de sentimentos e reconhecimento de entidades \cite{young_recent_2018, otter_survey_2021}.

Um marco decisivo nessa evolução foi a introdução da arquitetura \textit{Transformer}, proposta por \citeonline{vaswani_attention_2023}, que eliminou a necessidade de recorrências ao utilizar mecanismos de autoatenção. Com isso, foi possível o treinamento de modelos maiores e mais eficientes, conhecidos como \textit{Large Language Models} (LLMs). Esses modelos são capazes de gerar textos coerentes e contextualizados, demonstrando versatilidade inédita em múltiplas tarefas linguísticas \cite{chen_evaluating_2021}.

Os avanços proporcionados por LLMs e PLN vêm sendo amplamente aplicados para melhorar processos em diversas áreas. Na saúde, essas tecnologias têm auxiliado diagnósticos e agilizado a análise de dados clínicos \cite{paixao_machine_2022}. No setor jurídico, LLMs ajudam na pesquisa legal e previsão de resultados judiciais, tornando o trabalho mais eficiente \cite{lai_large_2024}. Atendimento ao cliente, administração pública e setor empresarial também têm integrado esses modelos para automatizar tarefas e aprimorar a tomada de decisão \cite{pessanha_large_2024}. Esse cenário abre possibilidades promissoras também para o campo educacional \cite{shahzad_comprehensive_2025}.

O sistema educacional brasileiro enfrenta desafios multidimensionais, abrangendo questões estruturais, pedagógicas, avaliativas e administrativas. Estruturalmente, persistem deficiências em infraestrutura e recursos limitados, especialmente na rede pública, problemas agravados por desigualdades socioeconômicas e regionais \cite{carnaval_desigualdade_2021}. Pedagogicamente, destaca-se a formação insuficiente de professores e práticas de ensino ultrapassadas, que afetam diretamente a qualidade do ensino \cite{gatti_formacao_2017}. No contexto administrativo, políticas públicas descontínuas, baixa valorização docente e deficiências no planejamento institucional impedem avanços consistentes \cite{silva_politicas_2020}.

Em relação às avaliações, prevalecem exames tradicionais centrados na memorização, negligenciando abordagens formativas e críticas. Essa realidade, aliada à gestão escolar deficiente, compromete a equidade e a qualidade das avaliações \cite{santos_articulacao_2016}.

Um desafio central nesse cenário educacional refere-se especificamente à correção de provas discursivas. Essa atividade exige alto empenho docente, análise minuciosa e critérios definidos para justiça na atribuição das notas. Estudos apontam que correções discursivas sofrem frequentemente de subjetividade e variação significativa entre avaliadores, impactando negativamente a equidade e a confiabilidade da avaliação educacional \cite{borges_avaliacao_2017}.


\section{Objetivos}

\subsection{Objetivo Geral} 
Desenvolver um sistema de apoio à correção de provas discursivas utilizando modelos avançados de linguagem (LLMs), apoiado pela recuperação dinâmica de contexto, de modo a produzir avaliações preliminares consistentes e alinhadas aos critérios pedagógicos estabelecidos pelos professores, com revisão e ajuste docente no fluxo final.

\subsection{Objetivos Específicos}
\begin{enumerate}
\item Construir um banco de dados vetorial a partir dos materiais didáticos fornecidos pelos professores, permitindo ao sistema recuperar automaticamente conteúdos relevantes durante o processo de avaliação das respostas discursivas. 

\item Desenvolver um módulo para configuração inicial das provas, em que os professores possam selecionar critérios avaliativos pré-existentes e definir, para cada critério, a pontuação máxima.

\item Implementar um mecanismo de correção baseado em múltiplos agentes virtuais independentes, capazes de analisar as respostas individualmente, tratando automaticamente eventuais divergências entre suas avaliações.

\item Disponibilizar uma interface intuitiva que permita aos professores revisar, ajustar e confirmar as correções realizadas pelo sistema, garantindo que as avaliações finais estejam plenamente alinhadas aos objetivos acadêmicos e critérios pedagógicos estabelecidos previamente.

\item Gerar análises estatísticas e de desempenho da avaliação, incluindo distribuição de notas, lacunas de aprendizagem por critério e tendência de desempenho por turma, de modo a apoiar a adaptação do planejamento de aulas.

\end{enumerate}

\section{Justificativa}

A correção de provas discursivas impõe uma demanda considerável sobre o tempo dos docentes, uma tarefa essencial que, no entanto, compete com outras atividades pedagógicas cruciais como o planejamento de aulas e o acompanhamento individualizado dos estudantes. Pesquisas sobre a jornada de trabalho docente indicam que uma parcela significativa de professores dedica entre 6 a 10 horas semanais a atividades extraclasse, que incluem a elaboração e correção de avaliações \cite{barbosa_tempo_2021}. Este cenário de sobrecarga não apenas afeta o bem-estar docente, mas possui o potencial de impactar a qualidade geral do ensino.

 A principal contribuição neste trabalho reside em oferecer ao professor uma ferramenta de suporte tecnológico capaz de agilizar a correção, liberando tempo para que o docente se dedique a aspectos mais estratégicos e humanizados da prática pedagógica. É crucial enfatizar que a solução proposta visa complementar, e não substituir, o papel central do professor, atuando como um assistente na análise inicial das respostas, sempre com base nos critérios definidos pelo educador.

A otimização do tempo docente é um benefício direto esperado. A literatura científica aponta que a IA no ensino superior pode ser um recurso valioso para personalizar o aprendizado e apoiar decisões pedagógicas, transformando a realidade dos alunos e também desafiando professores a se manterem atualizados sobre as tendências e sua integração ao fazer pedagógico, sem, contudo, substituir a mediação humana essencial ao processo de aprendizagem \cite{bottentuit_junior_inteligencia_2024} . Ao automatizar parcialmente a verificação da conformidade das respostas, espera-se que os professores possam reinvestir o tempo economizado no aprimoramento de suas práticas de ensino e na interação com os alunos, promovendo uma educação de maior qualidade \cite{gatti_formacao_2017}.

Adicionalmente, a correção manual está inerentemente sujeita à subjetividade, o que pode gerar inconsistências avaliativas \cite{borges_avaliacao_2017}. O sistema proposto busca oferecer uma avaliação inicial mais padronizada, embora a revisão final e os ajustes permaneçam sob a responsabilidade do professor, assegurando a conformidade com os objetivos acadêmicos e a experiência docente. Esta colaboração visa aumentar a equidade do processo avaliativo.

Em suma, a implementação desta solução busca contribuir para a melhoria da qualidade do ensino ao reduzir a carga operacional dos professores, permitindo que se concentrem em atividades que promovem a aprendizagem significativa. A ferramenta pode ainda auxiliar na identificação de padrões de dificuldades dos alunos, informando intervenções pedagógicas mais eficazes \cite{santos_articulacao_2016}. A proposta alinha-se a uma visão de que a tecnologia pode aprimorar os processos educativos, valorizando sempre o papel central do educador.

%\section{Estrutura do Trabalho}
